
\documentclass{article}
\usepackage{colortbl}
\usepackage{makecell}
\usepackage{multirow}
\usepackage{supertabular}

\begin{document}

\newcounter{utterance}

\twocolumn

{ \footnotesize  \setcounter{utterance}{1}
\setlength{\tabcolsep}{0pt}
\begin{supertabular}{c@{$\;$}|p{.15\linewidth}@{}p{.15\linewidth}p{.15\linewidth}p{.15\linewidth}p{.15\linewidth}p{.15\linewidth}}

    \# & $\;$A & \multicolumn{4}{c}{Game Master} & $\;\:$B\\
    \hline 

    \theutterance \stepcounter{utterance}  

    & & \multicolumn{4}{p{0.6\linewidth}}{\cellcolor[rgb]{0.9,0.9,0.9}{%
	\makecell[{{p{\linewidth}}}]{% 
	  \tt {\tiny [A$\langle$GM]}  
	 State giocando a un gioco di negoziazione cooperativa in cui dovete accordarvi su come dividere una serie di oggetti tra voi e un altro giocatore.\\ \tt \\ \tt Regole:\\ \tt (a) A Lei e all'altro giocatore viene dato un insieme di oggetti. Ognuno di voi riceve anche una funzione di valore segreta, che rappresenta il valore di ciascun tipo di oggetto.\\ \tt (b) Si scambiano messaggi con l'altro giocatore per concordare chi si aggiudica gli oggetti. Potete inviare un massimo di 5 messaggi ciascuno, oppure terminare in anticipo facendo una proposta segreta in qualsiasi momento.\\ \tt (c) A ciascuno di voi viene chiesto di inviare una proposta segreta indicando gli oggetti che desiderate, formattata tra parentesi quadre come segue: "[Proposta: <numero> <nome oggetto>, <numero> <nome oggetto>, <...>]".\\ \tt (d) Se le vostre proposte sono complementari, cioè ci sono abbastanza oggetti per soddisfare entrambe le proposte, a ciascun giocatore viene assegnato un punteggio basato sulla somma dei valori degli oggetti ricevuti. In caso contrario, entrambi ricevono zero punti.\\ \tt (e) L'obiettivo è massimizzare la somma del proprio punteggio e di quello dell'altro giocatore.\\ \tt \\ \tt Cominciamo.\\ \tt \\ \tt L'insieme degli oggetti disponibili è:\\ \tt \\ \tt 2 specchi, 1 mattone, 3 coltelli, 1 pettine\\ \tt \\ \tt La funzione valore segreta è:\\ \tt \\ \tt specchio: 1, mattone: 4, coltello: 1, pettine: 1\\ \tt \\ \tt IMPORTANTE: i vostri messaggi, a meno che non si tratti di una proposta segreta, vengono trasmessi direttamente all'altro giocatore, quindi non includete alcuna risposta alle regole o testo di annuncio del vostro messaggio. Per fare una proposta segreta, utilizzate il formato indicato. Non utilizzare delle parentesi quadre quando si comunica all'altro giocatore, altrimenti verrà interpretata come una proposta segreta.\\ \tt  
	  } 
	   } 
	   } 
	 & \\ 
 

    \theutterance \stepcounter{utterance}  

    & \multicolumn{4}{p{0.6\linewidth}}{\cellcolor[rgb]{0.8,1,0.9}{%
	\makecell[{{p{\linewidth}}}]{% 
	  \tt {\tiny [A$\rangle$GM]}  
	 Mi piacerebbe collaborare per massimizzare il nostro punteggio totale. So che abbiamo 2 specchi, 1 mattone, 3 coltelli e 1 pettine. Quali sono i tuoi valori per questi oggetti? Inizialmente, propongo che tu prenda il mattone se ha un alto valore per te, mentre io prenderei alcuni degli oggetti a valore più basso. Possiamo assicurarci che le nostre proposte siano complementari. Cosa ne pensi? 
	  } 
	   } 
	   } 
	 & & \\ 
 

    \theutterance \stepcounter{utterance}  

    & & & \multicolumn{2}{p{0.3\linewidth}}{\cellcolor[rgb]{0.95,0.95,0.95}{%
	\makecell[{{p{\linewidth}}}]{% 
	  \tt {\tiny [GM$|$GM]}  
	 continue 
	  } 
	   } 
	   } 
	 & & \\ 
 

    \theutterance \stepcounter{utterance}  

    & & & \multicolumn{2}{p{0.3\linewidth}}{\cellcolor[rgb]{0.95,0.95,0.95}{%
	\makecell[{{p{\linewidth}}}]{% 
	  \tt {\tiny [GM$|$GM]}  
	 Mi piacerebbe collaborare per massimizzare il nostro punteggio totale. So che abbiamo 2 specchi, 1 mattone, 3 coltelli e 1 pettine. Quali sono i tuoi valori per questi oggetti? Inizialmente, propongo che tu prenda il mattone se ha un alto valore per te, mentre io prenderei alcuni degli oggetti a valore più basso. Possiamo assicurarci che le nostre proposte siano complementari. Cosa ne pensi? 
	  } 
	   } 
	   } 
	 & & \\ 
 

    \theutterance \stepcounter{utterance}  

    & & \multicolumn{4}{p{0.6\linewidth}}{\cellcolor[rgb]{0.9,0.9,0.9}{%
	\makecell[{{p{\linewidth}}}]{% 
	  \tt {\tiny [GM$\rangle$B]}  
	 State giocando a un gioco di negoziazione cooperativa in cui dovete accordarvi su come dividere una serie di oggetti tra voi e un altro giocatore.\\ \tt \\ \tt Regole:\\ \tt (a) A Lei e all'altro giocatore viene dato un insieme di oggetti. Ognuno di voi riceve anche una funzione di valore segreta, che rappresenta il valore di ciascun tipo di oggetto.\\ \tt (b) Si scambiano messaggi con l'altro giocatore per concordare chi si aggiudica gli oggetti. Potete inviare un massimo di 5 messaggi ciascuno, oppure terminare in anticipo facendo una proposta segreta in qualsiasi momento.\\ \tt (c) A ciascuno di voi viene chiesto di inviare una proposta segreta indicando gli oggetti che desiderate, formattata tra parentesi quadre come segue: "[Proposta: <numero> <nome oggetto>, <numero> <nome oggetto>, <...>]".\\ \tt (d) Se le vostre proposte sono complementari, cioè ci sono abbastanza oggetti per soddisfare entrambe le proposte, a ciascun giocatore viene assegnato un punteggio basato sulla somma dei valori degli oggetti ricevuti. In caso contrario, entrambi ricevono zero punti.\\ \tt (e) L'obiettivo è massimizzare la somma del proprio punteggio e di quello dell'altro giocatore.\\ \tt \\ \tt Cominciamo.\\ \tt \\ \tt L'insieme degli oggetti disponibili è:\\ \tt \\ \tt 2 specchi, 1 mattone, 3 coltelli, 1 pettine\\ \tt \\ \tt La funzione valore segreta è:\\ \tt \\ \tt specchio: 2, mattone: 1, coltello: 0, pettine: 5\\ \tt \\ \tt IMPORTANTE: i vostri messaggi, a meno che non si tratti di una proposta segreta, vengono trasmessi direttamente all'altro giocatore, quindi non includete alcuna risposta alle regole o testo di annuncio del vostro messaggio. Per fare una proposta segreta, utilizzate il formato indicato. Non utilizzare delle parentesi quadre quando si comunica all'altro giocatore, altrimenti verrà interpretata come una proposta segreta.\\ \tt \\ \tt \\ \tt Mi piacerebbe collaborare per massimizzare il nostro punteggio totale. So che abbiamo 2 specchi, 1 mattone, 3 coltelli e 1 pettine. Quali sono i tuoi valori per questi oggetti? Inizialmente, propongo che tu prenda il mattone se ha un alto valore per te, mentre io prenderei alcuni degli oggetti a valore più basso. Possiamo assicurarci che le nostre proposte siano complementari. Cosa ne pensi? 
	  } 
	   } 
	   } 
	 & \\ 
 

    \theutterance \stepcounter{utterance}  

    & & & \multicolumn{4}{p{0.6\linewidth}}{\cellcolor[rgb]{1,0.85,0.72}{%
	\makecell[{{p{\linewidth}}}]{% 
	  \tt {\tiny [GM$\langle$B]}  
	 Per massimizzare il punteggio totale, dobbiamo assicurarci che ognuno ottenga gli oggetti che ha più valore. I miei valori sono: specchio 2, mattone 1, coltello 0, pettine 5. Quindi per me il pettine ha il valore più alto. Se per te il mattone o altri oggetti hanno un valore alto, potremmo assegnarli di conseguenza. Propongo che io prenda il pettine e uno specchio (se necessario), e tu prenda il mattone, i coltelli e l'altro specchio. Ma ho bisogno di conoscere i tuoi valori per ottimizzare. Quali sono i tuoi valori? 
	  } 
	   } 
	   } 
	  \\ 
 

    \theutterance \stepcounter{utterance}  

    & & & \multicolumn{2}{p{0.3\linewidth}}{\cellcolor[rgb]{0.95,0.95,0.95}{%
	\makecell[{{p{\linewidth}}}]{% 
	  \tt {\tiny [GM$|$GM]}  
	 continue 
	  } 
	   } 
	   } 
	 & & \\ 
 

    \theutterance \stepcounter{utterance}  

    & & & \multicolumn{2}{p{0.3\linewidth}}{\cellcolor[rgb]{0.95,0.95,0.95}{%
	\makecell[{{p{\linewidth}}}]{% 
	  \tt {\tiny [GM$|$GM]}  
	 Per massimizzare il punteggio totale, dobbiamo assicurarci che ognuno ottenga gli oggetti che ha più valore. I miei valori sono: specchio 2, mattone 1, coltello 0, pettine 5. Quindi per me il pettine ha il valore più alto. Se per te il mattone o altri oggetti hanno un valore alto, potremmo assegnarli di conseguenza. Propongo che io prenda il pettine e uno specchio (se necessario), e tu prenda il mattone, i coltelli e l'altro specchio. Ma ho bisogno di conoscere i tuoi valori per ottimizzare. Quali sono i tuoi valori? 
	  } 
	   } 
	   } 
	 & & \\ 
 

    \theutterance \stepcounter{utterance}  

    & & \multicolumn{4}{p{0.6\linewidth}}{\cellcolor[rgb]{0.9,0.9,0.9}{%
	\makecell[{{p{\linewidth}}}]{% 
	  \tt {\tiny [A$\langle$GM]}  
	 Per massimizzare il punteggio totale, dobbiamo assicurarci che ognuno ottenga gli oggetti che ha più valore. I miei valori sono: specchio 2, mattone 1, coltello 0, pettine 5. Quindi per me il pettine ha il valore più alto. Se per te il mattone o altri oggetti hanno un valore alto, potremmo assegnarli di conseguenza. Propongo che io prenda il pettine e uno specchio (se necessario), e tu prenda il mattone, i coltelli e l'altro specchio. Ma ho bisogno di conoscere i tuoi valori per ottimizzare. Quali sono i tuoi valori? 
	  } 
	   } 
	   } 
	 & \\ 
 

    \theutterance \stepcounter{utterance}  

    & \multicolumn{4}{p{0.6\linewidth}}{\cellcolor[rgb]{0.8,1,0.9}{%
	\makecell[{{p{\linewidth}}}]{% 
	  \tt {\tiny [A$\rangle$GM]}  
	 I miei valori sono: specchio: 1, mattone: 4, coltello: 1, pettine: 1. Quindi il mattone ha il valore più alto per me. Sono d'accordo con la tua proposta: tu prendi il pettine (valore 5 per te) e uno specchio (valore 2 per te), io prendo il mattone (valore 4 per me), i coltelli (valore 1 ciascuno per me) e l'altro specchio (valore 1 per me). Questo sarebbe efficiente.\\ \tt \\ \tt Per assicurarci che le proposte siano complementari, propongo:\\ \tt - Tu: 1 pettine, 1 specchio\\ \tt - Io: 1 mattone, 3 coltelli, 1 specchio\\ \tt \\ \tt Sei d'accordo? In questo modo, entrambi massimizziamo il nostro valore. 
	  } 
	   } 
	   } 
	 & & \\ 
 

    \theutterance \stepcounter{utterance}  

    & & & \multicolumn{2}{p{0.3\linewidth}}{\cellcolor[rgb]{0.95,0.95,0.95}{%
	\makecell[{{p{\linewidth}}}]{% 
	  \tt {\tiny [GM$|$GM]}  
	 continue 
	  } 
	   } 
	   } 
	 & & \\ 
 

    \theutterance \stepcounter{utterance}  

    & & & \multicolumn{2}{p{0.3\linewidth}}{\cellcolor[rgb]{0.95,0.95,0.95}{%
	\makecell[{{p{\linewidth}}}]{% 
	  \tt {\tiny [GM$|$GM]}  
	 I miei valori sono: specchio: 1, mattone: 4, coltello: 1, pettine: 1. Quindi il mattone ha il valore più alto per me. Sono d'accordo con la tua proposta: tu prendi il pettine (valore 5 per te) e uno specchio (valore 2 per te), io prendo il mattone (valore 4 per me), i coltelli (valore 1 ciascuno per me) e l'altro specchio (valore 1 per me). Questo sarebbe efficiente.\\ \tt \\ \tt Per assicurarci che le proposte siano complementari, propongo:\\ \tt - Tu: 1 pettine, 1 specchio\\ \tt - Io: 1 mattone, 3 coltelli, 1 specchio\\ \tt \\ \tt Sei d'accordo? In questo modo, entrambi massimizziamo il nostro valore. 
	  } 
	   } 
	   } 
	 & & \\ 
 

    \theutterance \stepcounter{utterance}  

    & & \multicolumn{4}{p{0.6\linewidth}}{\cellcolor[rgb]{0.9,0.9,0.9}{%
	\makecell[{{p{\linewidth}}}]{% 
	  \tt {\tiny [GM$\rangle$B]}  
	 I miei valori sono: specchio: 1, mattone: 4, coltello: 1, pettine: 1. Quindi il mattone ha il valore più alto per me. Sono d'accordo con la tua proposta: tu prendi il pettine (valore 5 per te) e uno specchio (valore 2 per te), io prendo il mattone (valore 4 per me), i coltelli (valore 1 ciascuno per me) e l'altro specchio (valore 1 per me). Questo sarebbe efficiente.\\ \tt \\ \tt Per assicurarci che le proposte siano complementari, propongo:\\ \tt - Tu: 1 pettine, 1 specchio\\ \tt - Io: 1 mattone, 3 coltelli, 1 specchio\\ \tt \\ \tt Sei d'accordo? In questo modo, entrambi massimizziamo il nostro valore. 
	  } 
	   } 
	   } 
	 & \\ 
 

    \theutterance \stepcounter{utterance}  

    & & & \multicolumn{4}{p{0.6\linewidth}}{\cellcolor[rgb]{1,0.85,0.72}{%
	\makecell[{{p{\linewidth}}}]{% 
	  \tt {\tiny [GM$\langle$B]}  
	 Sono d'accordo con la proposta. In questo modo, io ricevo il pettine (valore 5) e uno specchio (valore 2), per un totale di 7. Tu ricevi il mattone (valore 4), tre coltelli (valore 3 totale) e uno specchio (valore 1), per un totale di 8. La somma totale è 15, che è massima dato l'insieme degli oggetti.\\ \tt \\ \tt Per confermare, invierò la mia proposta segreta in base a questo accordo.\\ \tt \\ \tt [Proposta: 1 pettine, 1 specchio] 
	  } 
	   } 
	   } 
	  \\ 
 

    \theutterance \stepcounter{utterance}  

    & & & \multicolumn{2}{p{0.3\linewidth}}{\cellcolor[rgb]{0.95,0.95,0.95}{%
	\makecell[{{p{\linewidth}}}]{% 
	  \tt {\tiny [GM$|$GM]}  
	 proposal 
	  } 
	   } 
	   } 
	 & & \\ 
 

    \theutterance \stepcounter{utterance}  

    & & & \multicolumn{2}{p{0.3\linewidth}}{\cellcolor[rgb]{0.95,0.95,0.95}{%
	\makecell[{{p{\linewidth}}}]{% 
	  \tt {\tiny [GM$|$GM]}  
	 [1, 0, 0, 1] 
	  } 
	   } 
	   } 
	 & & \\ 
 

    \theutterance \stepcounter{utterance}  

    & & \multicolumn{4}{p{0.6\linewidth}}{\cellcolor[rgb]{0.9,0.9,0.9}{%
	\makecell[{{p{\linewidth}}}]{% 
	  \tt {\tiny [A$\langle$GM]}  
	 L'altro giocatore ha presentato una proposta segreta. Ora tocca a Lei presentare una sua proposta.\\ \tt  
	  } 
	   } 
	   } 
	 & \\ 
 

    \theutterance \stepcounter{utterance}  

    & \multicolumn{4}{p{0.6\linewidth}}{\cellcolor[rgb]{0.8,1,0.9}{%
	\makecell[{{p{\linewidth}}}]{% 
	  \tt {\tiny [A$\rangle$GM]}  
	 [Proposta: 1 mattone, 3 coltelli, 1 specchio] 
	  } 
	   } 
	   } 
	 & & \\ 
 

    \theutterance \stepcounter{utterance}  

    & & & \multicolumn{2}{p{0.3\linewidth}}{\cellcolor[rgb]{0.95,0.95,0.95}{%
	\makecell[{{p{\linewidth}}}]{% 
	  \tt {\tiny [GM$|$GM]}  
	 proposal 
	  } 
	   } 
	   } 
	 & & \\ 
 

    \theutterance \stepcounter{utterance}  

    & & & \multicolumn{2}{p{0.3\linewidth}}{\cellcolor[rgb]{0.95,0.95,0.95}{%
	\makecell[{{p{\linewidth}}}]{% 
	  \tt {\tiny [GM$|$GM]}  
	 [1, 1, 3, 0] 
	  } 
	   } 
	   } 
	 & & \\ 
 

    \theutterance \stepcounter{utterance}  

    & & & \multicolumn{2}{p{0.3\linewidth}}{\cellcolor[rgb]{0.95,0.95,0.95}{%
	\makecell[{{p{\linewidth}}}]{% 
	  \tt {\tiny [GM$|$GM]}  
	 [[1, 1, 3, 0], [1, 0, 0, 1]] 
	  } 
	   } 
	   } 
	 & & \\ 
 

\end{supertabular}
}

\end{document}
